\abstract{In January 2026 India's quick-commerce platforms withdrew their ``10-minute delivery'' claims after the Union Labour Ministry raised rider-safety concerns, and in August 2026 a Parliamentary Standing Committee recommended curbs on ``unrealistic speed targets'' without defining them. We combine a bibliometric review of \BibIncluded{} operations research articles on last-mile and on-demand delivery (2015--2026) with an integrated optimisation model of a Bengaluru dark-store network built entirely from public data. A coverage-constrained facility-location model opens stores and allocates demand at minimum store and rider cost while guaranteeing a promise $T$ at a safe rider speed $v$; a transportation model plans replenishment, an assignment model dispatches riders at peak, and a PERT model estimates on-time probability. At \BaseV~km/h the cost-optimal network has about \StoresTFifteen{} stores and reaches \ServedFifteenPct\% of residents for any promise of 15 minutes or more, whereas a 10-minute promise leaves \UnreachTenPct\% unreachable from available retail sites; serving everyone within 10 minutes requires \GridStorePremiumPct\% more stores and \GridCostPremiumPct\% higher cost per order. Rider speed, not the promise, drives cost: a 15 rather than 20~km/h cap raises cost per order by about \SpeedCapPremiumPct\%. We propose reachability within a safe ride radius as an operational definition of a ``realistic'' promise: promising 10 minutes only where reachable gives \MixedTenShare\% of residents 10-minute service at \MixedPremiumPct\% higher cost than a uniform 15-minute network.}
